\section{Evaluation Strategy}

In terms of evaluation, a core comparison will be performed for each architecture. This will involve
comparing the base model vs SFT vs OPSD. The same train, validation and test split will be used
for each method/model used. Performance will be measured using BERTScore-F1 as the primary metric
together with normalised exact-match accuracy as a secondary metric.

BERTScore compares generated and ground-truth answers using contextual token embeddings making it
suitable for open-ended VQA where semantically equivalent answers are still correct. Normalised exact-match 
accuracy will provide a stricter complementary metric. To ensure this metric is used correctly,
outputs will be preprocessed for case, punctuation, articles and formatting differences.

For RQ1, the base model, SFT and OPSD performances will be compared within each model. This will determine
whether OPSD provides a consistent performance improvement over SFT and the base model without fine-tuning. For RQ2,
the magnitude of this improvement will be compared across the different models tested. This enables one to 
assess whether the effectiveness of OPSD varies according to model size and architecture.

For this experiment, there are two main baselines. The base model serves as the pre-fine-tuned baseline.
This establishes the performance of a pretrained VLM on the Visual Genome dataset before any task-specific
adaptation. The main methodological baseline is the SFT model. This is the main comparator for OPSD. The OPSD model
is the proposed method.